\section{Choosing the right plane description}

A plane can be described in several ways. The best description depends on what information is given and what question is being asked.

This is the same lesson that appeared for lines. A formula is not the object itself. A formula is a language for reading the object.

\subsection*{When a normal vector is given}

If a point and a normal vector are given, use the normal form:
\[
n\cdot(P-P_0)=0.
\]
In coordinates, if \(n=(A,B,C)\), this becomes
\[
A(x-x_0)+B(y-y_0)+C(z-z_0)=0.
\]
This form is usually the fastest route to the general equation.

\subsection*{When three points are given}

If three non-collinear points are given, first build two direction vectors:
\[
B-A,
\qquad
C-A.
\]
Then choose whether to write the plane parametrically or compute a normal vector using the cross product:
\[
n=(B-A)\times(C-A).
\]
The normal vector then gives the general equation.

\subsection*{When an equation is given}

If a plane is given by
\[
Ax+By+Cz+D=0,
\]
then the normal vector is immediately visible:
\[
n=(A,B,C).
\]
This is useful for distances, angles, parallelism, and perpendicularity.

\subsection*{Which form helps which question?}

\begin{center}
\begin{tabular}{ll}
\toprule
Question & Useful description \\
\midrule
Plane through point and normal vector & normal form \\
Plane through three points & two direction vectors and cross product \\
Generate points on the plane & parametric form \\
Check whether a point lies on the plane & general equation or parametric form \\
Find distance from point to plane & general equation \\
Angle between two planes & normal vectors \\
Parallel or perpendicular planes & normal vectors \\
Line-plane intersection & substitute line into plane equation \\
\bottomrule
\end{tabular}
\end{center}

The table should be read as a guide. A good solution often converts from one description to another.

\subsection*{A complete reading example}

Consider the plane
\[
2x-y+z-4=0.
\]
From the equation we immediately read a normal vector:
\[
n=(2,-1,1).
\]
A point on the plane can be found by choosing \(y=0\) and \(z=0\). Then
\[
2x-4=0,
\]
so
\[
x=2.
\]
Thus \((2,0,0)\) lies on the plane.

To find two direction vectors, choose any two vectors perpendicular to \(n\). For example,
\[
u=(1,2,0)
\]
works because
\[
(2,-1,1)\cdot(1,2,0)=2-2+0=0.
\]
Also,
\[
v=(0,1,1)
\]
works because
\[
(2,-1,1)\cdot(0,1,1)=0-1+1=0.
\]
Therefore one parametric description is
\[
P=(2,0,0)+s(1,2,0)+t(0,1,1).
\]

The same plane has been read in two languages: the general equation shows the normal vector, while the parametric equation shows directions inside the plane.

\begin{summarybox}
At the end of this chapter, remember:
\begin{itemize}
\item A plane needs either two independent directions or one normal direction.
\item The equation \(Ax+By+Cz+D=0\) contains the normal vector \((A,B,C)\).
\item Three points determine a plane only when they are not collinear.
\item Distances and angles are easiest with normal vectors.
\item A line-plane relationship is studied by combining the line parametrization with the plane equation.
\end{itemize}
\end{summarybox}

The chapter completes the analytic geometry sequence: points, vectors, lines, and planes are now all described through coordinates and vector language.